% Working input fragment. This is a new arithmetic construction, not an exact
% equivalence of the complete old witness space. No compilation claim is made.
\section{A shared-point construction}
\label{sec:new-shared-source}

\paragraph{Scope.}
The construction below assumes the source manuscript's cyclic-group,
five-contact height, scalar-height, and pole-completeness results, and also
Theorem~4.1 of Sun, \emph{86 unknowns for Hilbert's tenth problem over the
rationals}, arXiv:2610.20691v2. The last theorem is an additional arithmetic
input. The number of variables is not inferred from Sun's undecidability
theorem. All algebraic operations below have rational coefficients; field
operations in $K=\mathbb Q(\sqrt2)$ mean their two rational components.

\subsection{Fixed data and the complete rational system}
Use
\[
 d=3-\sqrt2,\quad D=d^2=11-6\sqrt2,\quad
 E:y^2=x^3-Dx,
\]
\[
 P=(57/8-7\sqrt2/2,-165/8+417\sqrt2/32),\qquad
 2E(K)=\mathbb ZP.
\]
For a nonzero point put $f(Q)=1/x(Q)$ and $z(Q)=-x(Q)/y(Q)$.
Write $\sigma$ for conjugation. Every displayed quotient is accompanied
by the inequation that its denominator, or its field norm, is nonzero.
The old point charts are replaced by doubles of affine halfpoints.

The ten point occurrences are
\[
 A,S,T,Z,J_0,H,U_A,U_S,U_T,U_Z.
\]
Every occurrence is nonidentity and lies in $2E(K)$. The points $A,S,T,Z$
have integer labels $(n,a,t,b)$, respectively; these labels are used only
in the proof and are not additional rational coordinates.

Set $\beta=X^2+2$. Use rational proxies $\xi,\tau,\zeta$ with
\[
 \tau^2-\xi(\xi-2)\tau\zeta+(\xi-2)^2\zeta^2-1=0.
\]
Only $\xi,\zeta$ are retained; $\tau$ is a finite algebraic coordinate.
Let $Q_0=P$, $(Q_1,Q_2,Q_3,Q_4)=(A,S,T,Z)$ and
$(U_0,U_1,U_2,U_3,U_4)=(H,U_A,U_S,U_T,U_Z)$.
Form $V_i=U_i-[4]Q_i$ for $0\le i\le4$. These five differences are
required to be nonidentity. Affine subtraction also retains the guard
$x(U_i)\ne x([4]Q_i)$; completeness will choose a common index different
from $0,4,-4$.

For $e\in\{16,64\}$ put
\[
 L_e(Q)=z(Q)G_{e/4-1}(Df(Q)^2),\quad
 G_m(W)=1+\sum_{j=1}^m c_jW^j,
\]
where $c_1=-2/5$ and $c_j=c_{j-1}(4j-5)/(4j+1)$.
The source proves these logarithms nonzero at every required nonidentity
point. For $k=(16,64,64,64)$ and $b=(\beta,\xi,\tau,\zeta)$ define
\[
 L_i=\frac{L_{k_i}(U_i)}{L_{64}(H)},\qquad
 L'_i=\sigma\!\left(\frac{L_{k_i}(V_i)}{L_{64}(V_0)}\right).
\]
The eighteen field tests $D_j$, with orders $m_j$, are
\[
 \begin{array}{c|c}
 f(U_0),\ldots,f(U_4),\sigma f(V_0),\ldots,\sigma f(V_4)&1\ (10\text{ entries})\\
 L_i-L'_i\quad(1\le i\le4)& k_i\\
 L_i-b_i\quad(1\le i\le4)&k_i.
 \end{array}
\]
There is no approximate square root of two and no common comparison
point $J$ in this construction.

Enumerate the eight rational components of $x(A),x(S),x(T),x(Z)$ as
$x_1,\ldots,x_8$. On finite branches set $A_0=1$ and, for $1\le r\le8$,
set $A_r=x_r^2$ if $x_r\ne0$, and $A_r=1$ if $x_r=0$. These are
expressions, not new variables. The common slope is one rational
coordinate $\lambda$. Put
\[
 \Theta_0=\{0,-1/2,-1,1/2399,-32/407\},\qquad
 F_5(T)=\prod_{\theta\in\Theta_0}(T-\theta),\quad t_r=A_r\lambda.
\]
Require $t_r\notin\Theta_0\cup\{1\}$, and define, for $0\le r\le8$,
\[
 R_r=A_r^{-5}\frac{F_5(t_r)}{(t_r-1)^6},\qquad
 \widetilde D_{r,j}=\frac{D_j}{(t_r-1)^{m_j+1}},
\]
\[
 W_r=\prod_{j=1}^{18}\left(1+R_r
       \left(\frac{\widetilde D_{r,j}}{R_r^{m_j}}\right)^2\right),
 \quad T_r=\operatorname{Tr}_{K/\mathbb Q}(W_r^{-1}),
 \quad B_r=R_r/T_r^2.
\]
Require $W_r,T_r\ne0$ and the following pole condition for each $B_r$.

Retain two rational coordinates $e_r,s_r$, set
\[
 C_r=17B_r,\qquad h_r=e_r(C_r-B_r^2),
\]
and impose the two monic equations
\[
 p_r^2=C_r(1-e_r)(1+3e_r),\qquad
 q_r^2=\frac{(s_r^2-C_rh_r)(s_r^2+3C_rh_r)}{C_r}.
\]
The roots $p_r,q_r$ are finite algebraic coordinates. This is precisely
the old pole predicate with preparation variable zero. No new finite
multiplier choice is needed. The prime-pattern theorem will be invoked
with its odd squarefree avoidance modulus divisible by $3\cdot17$.
It then makes each selected prime square in $\mathbb Q_2,\mathbb Q_3$
and $\mathbb Q_{17}$. Thus the fixed multiplier $17$ works throughout.

For use in the scalar-height block alone, retain the finite set
\[
 \mathcal C=\left\{\pm\prod_{p\in I}p:
 I\subseteq\{2,3,5,7,11,13,17\}\right\}.
\]

\paragraph{The unchanged scalar-height block.}
Put
\[
 r_\beta=\beta(\beta-1)(\beta-2401)(32\beta+343)(2\beta-1).
\]
Take the six rational coefficients of
$f(J_0),f(S),z(S)/z(J_0)$, with orders $(1,1,1,1,48,48)$.
Use the source's fixed binary norm construction $N_r(a,b)=a^2+rb^2$
on the constant leaf $1$ and the six normalized tests, and denote its
root by $W$. On a finite branch choose $c\in\mathcal C$ and set
$b_*=cr_\beta W^2$, $B_*=17b_*$. Retain $t_*,e_*,s_*$ and define
\[
 U_* =b_*\frac{1-b_*t_*^2}{1+b_*t_*^2},\qquad
 h_*=e_*(B_*-U_*^2).
\]
Require $b_*\ne0$, $1+b_*t_*^2\ne0$, and the two monic equations
\[
 p_*^2=B_*(1-e_*)(1+3e_*),\qquad
 q_*^2=\frac{(s_*^2-B_*h_*)(s_*^2+3B_*h_*)}{B_*}.
\]
The extra guard $b_*\ne0$ removes no required integral-input witnesses:
the source constructs $W$ as a unit at a prime dividing the nonzero
integer $r_\beta$, which has absolute value greater than one.

\paragraph{A fixed finite size condition.}
Let $H$ be the original five-contact constant, with
$1\le H<2^{4783227}$, and put $M=2^{9566463}$. Require
\[
 \mathcal A_M(S):=\prod_{j=1}^M
 N_{K/\mathbb Q}\bigl(x(S)-x([j]P)\bigr)\ne0.
\]
This is an indexed finite expression with fixed rational coefficients.
It forces $|a|>M>(20H)^2$. It has not been expanded. An alternative
source replaces this guard by the proved three-witness semilocal anchor;
that alternative adds two retained coordinates and one quadratic root.

\subsection{The finite algebra and the variable count}
For a point occurrence retain $a_0,a_1$ and set $a=a_0+a_1\sqrt2$,
$a^3-Da=A+B\sqrt2$. Introduce a root $C$ of the monic quartic
\[
 C^4-AC^2+B^2/2=0,
\]
require $C\ne0$ and $A^2-2B^2\ne0$, and put
\[
 b=C+\frac{B}{2C}\sqrt2.
\]
The quartic gives $b^2=a^3-Da$. Use $[2](a,b)$ for the point,
retaining all necessary norm guards. This single chart suffices for
the specific curve used here; there is no finite chart parameter.

\begin{lemma}[A single ordinate chart over a quadratic field]
Let $F$ have characteristic different from two, let
$K=F(\sqrt\delta)$ be a quadratic field, and let $d\in K^\times$
satisfy $d^3\notin F\cup\sqrt\delta F$. On
$E_d:y^2=x^3-d^2x$, every nonidentity $Q\in2E_d(K)$ has a
halfpoint whose ordinate has nonzero $F$-component. In particular
this applies to our curve. More generally it applies over a real
quadratic field whenever $d=u+v\sqrt\delta$ with $uv\ne0$.
\end{lemma}
\begin{proof}
Start with any affine halfpoint, so $b\ne0$ and $a\ne0$.
Translation by $T_0=(0,0)$ gives
\[
 (a,b)+T_0=(-d^2/a,d^2b/a^2)
\]
and does not change its double. If both ordinates had $F$-component
zero, then
$b$ and $d^2b/a^2$ would be nonzero $F$-multiples of $\sqrt\delta$,
so $d^2/a^2\in F$. Thus $w=a/d$ has square in $F$ and belongs to
$F$ or $\sqrt\delta F$. But $b^2=d^3w(w^2-1)\in F^\times$.
In the first case this forces $d^3\in F$, and in the second it
forces $d^3\in\sqrt\delta F$, both contradictions. The nonzero
factor $w(w^2-1)$ follows from $b\ne0$.

Over a real quadratic field, the two components of
$d^3$ are $u(u^2+3\delta v^2)$ and $v(3u^2+\delta v^2)$,
both nonzero when $uv\ne0$. For the present curve, explicitly,
\[
 d^3=45-29\sqrt2,\qquad
 d^3\sqrt2=-58+45\sqrt2.
\]
The addition formula has the displayed positive sign: its slope is
$b/a$, its new abscissa is $-d^2/a$, and its new ordinate is
$(b/a)(a+d^2/a)-b=d^2b/a^2$.
\end{proof}

The quartic derivative is nonzero at the permitted roots. Indeed,
with $Y=C^2$, the equation is $Y^2-AY+B^2/2=0$, of discriminant
$A^2-2B^2\ne0$, and $C\ne0$. Thus its derivative may be recorded
explicitly as a guard without excluding any permitted root.

The retained coordinates number
\[
 1+20+2+1+18+3=45.
\]
The remaining coordinates are ten quartic roots, eighteen pole roots,
two scalar-height roots and the Lucas root. Their monic triangular
degree product is $4^{10}2^{21}=2^{41}$. Vanishing square radicands
are treated by the linear relation that their root is zero; on each
nonzero branch its root is guarded nonzero. The finite-fibre theorem
therefore supplies one additional rational coordinate, for a total of
46. Finite choices and their strata are combined by products, not by
new existential selectors.

\subsection{Local tests and the fixed preparation-zero predicate}

\begin{lemma}[A restricted completeness domain with preparation zero]
The preparation-zero pole just specified excludes negative odd
valuations outside $\{2,3,17\}$. It has rational witnesses whenever
$b$ has square reduced denominator, is positive and square in
$\mathbb Q_2,\mathbb Q_3,\mathbb Q_{17}$, and satisfies
\[
 b<17,\qquad v_2(b)\ge3,\quad v_3(b)\ge1,\quad v_{17}(b)\ge2.
\]
This assertion uses only the pole-completeness arithmetic input already
used by the source.
\end{lemma}
\begin{proof}
The equations are the old equations with preparation variable zero and
$U=b$, so the old one-sided valuation obstruction applies. For
completeness put $B=17b$ and $V=B(B-b^2)=17^2b^2(1-b/17)$.
The stated inequalities make $V$ positive and square at the three
specified local fields. They are exactly the sufficient preparation
conditions in the source's Lemma
\texttt{original-seven-parity-completeness}. Its remaining argument
constructs rational points in square class $B$ on the two twists and
then rescales them to the displayed equations. No further preparation
variable is used. Equivalently one may repeat that argument with
$U=b$: at a support prime of the square class of $B$ outside the three
exceptions, $v_p(b)$ is positive odd, so $1-b/17$ is a square unit;
at a support prime of $V$ the valuation of $B$ is even and cancellation
forces $B/b^2$ to be a square unit. Thus the same disjoint-support and
mutual-squareness hypotheses for the simultaneous two-twist construction
hold. Its conclusions supply the required witnesses.
\end{proof}

\begin{lemma}[One place for all local tests]
Let $R\ne0$ be rational, let $D_j\in K$, and let $m_j\ge1$.
Set
\[
 W=\prod_j(1+R(D_j/R^{m_j})^2),\quad
 T=\operatorname{Tr}_{K/\mathbb Q}(W^{-1}),\quad b=R/T^2,
\]
with the displayed nonzero guards and a valid preparation-zero pole
certificate for $b$. If $p>76176$ and $v_p(R)$ is odd, then it is
positive and some place $\mathfrak p\mid p$ satisfies
$v_{\mathfrak p}(D_j)\ge m_j$ for every $j$.
\end{lemma}
\begin{proof}
Put $\alpha=v_p(R)$. All valuations in $K_{\mathfrak p}$ are integral,
because $p$ is unramified. Each nonzero nonconstant term of a factor
has odd valuation
$b_{j,\mathfrak p}=\alpha+2(v_{\mathfrak p}(D_j)-m_j\alpha)$,
so its factor has valuation $\min(0,b_{j,\mathfrak p})$. Put
$\lambda_{\mathfrak p}=v_{\mathfrak p}(W)\le0$ and
$\lambda=\max_{\mathfrak p\mid p}\lambda_{\mathfrak p}\le0$.
Integrality of trace gives $v_p(T)\ge-\lambda$. Therefore
\[
 v_p(b)=\alpha-2v_p(T)\le\alpha+2\lambda.
\]
The left side is odd and cannot be negative by the pole obstruction.
It is consequently positive, so $\alpha>0$ and
$\lambda>-\alpha/2$. At a place attaining $\lambda$, each nonpositive
summand of $v(W)$ is at least $\lambda$. For each nonzero $D_j$ this gives
\[
 \alpha+2(v(D_j)-m_j\alpha)>-\alpha/2,
 \qquad v(D_j)>(m_j-3/4)\alpha\ge m_j-3/4.
\]
Integral valuations yield $v(D_j)\ge m_j$. Zero terms meet the same
conclusion automatically. This is the reciprocal-trace argument of
Sun's Lemma~3.1, written here for one shared group.
\end{proof}

\begin{lemma}[Two anchors compare the four integer labels]
At a prime $p>76176$, suppose the eighteen field tests have their stated
orders at a place $\mathfrak p\mid p$. Then every proxy $b_i$ is
$p$-integral and
\[
 b_i\equiv n_i\pmod{p^{k_i}},\qquad
 (n_1,n_2,n_3,n_4)=(n,a,t,b).
\]
This statement does not assume that $p$ splits or that the label of
$H$ is a $p$-unit.
\end{lemma}
\begin{proof}
Write $H=[h]P$, $U_i=[u_i]P$. The first ten tests place $H,U_i$ in
the formal group at $\mathfrak p$, and $V_0,V_i$ in the formal group
at its conjugate place. The guarded denominators ensure $h\ne0,4$.
The source's formal logarithm identities give the exact ratios
\[
 r_i=u_i/h,\qquad r'_i=(u_i-4n_i)/(h-4).
\]
The reciprocal truncations have the same valuations as these ratios
and agree with them to relative precision $k_i$. At least one of $h$
and $h-4$ is a $p$-unit; its corresponding exact ratio is integral.
The test $v(L_i-L'_i)\ge k_i$ then makes the other ratio integral too.
Consequently the relative precision becomes absolute precision for
both ratios. The remaining test gives
$b_i\equiv r_i\equiv r'_i\pmod{p^{k_i}}$. Multiplying by the integer
denominators and subtracting yields
$4(b_i-n_i)\equiv0\pmod{p^{k_i}}$. Since $p\ne2$, the assertion follows.
If $p$ is inert, conjugation acts within the same local field and the
same proof applies. No splitting assertion is needed for soundness.
\end{proof}

\subsection{Height transfer with the original numerical constant}
Let $M_0(t)$ denote the odd-contact radical for the five new contacts,
with primes at most $76176$ omitted. The map $\phi(t)=2+1/t$ carries
these contacts to the original five contacts. Its matrix and inverse
are integral with determinant $-1$, and the coordinate forms differ
by fixed rational factors. The prime divisors of the new nonzero
coefficients and contact differences, including differences from $1$,
are
\[
 2,3,5,7,11,37,109,439,2399.
\]
Thus no new exceptional prime is introduced. For $t$ avoiding the
five contacts,
\[
 M_0(t)=M_5(\phi(t)),\qquad
 h(t)\le h(\phi(t))+\log3
       \le(H+\log3)M_0(t)^{24}.
\]

At a prime in $M_0(t_r)$, $t_r-1$ is a unit and exactly one factor
of $F_5(t_r)$ has positive odd valuation. Because $A_r^{-5}$ is a
rational square, $v_p(R_r)$ is odd even when $A_r$ is not a unit.
The preceding trace lemma applied to $\widetilde D_{r,j}$ therefore
gives the original, unscaled tests to their stated precision. Hence
the two-anchor lemma applies at every prime in every $M_0(t_r)$.

Suppose the $K$th powers of all nine radicals divide a nonzero integer
$N$. Then
\[
 h(\lambda),h(A_r\lambda)
 \le(H+\log3)|N|^{24/K}.
\]
For $x_r\ne0$, the identity $A_r=x_r^2$ gives
\[
 2h(x_r)=h(A_r)\le h(\lambda)+h(A_r\lambda).
\]
The same inequality for $h(x_r)$ is immediate at $x_r=0$. Finally,
\[
 h(x(Q_i))\le h(x_{2i-1})+h(x_{2i})+\tfrac32\log2
 <7H|N|^{24/K},
\]
because $H\ge1$, $|N|\ge1$, and $2\log3+\tfrac32\log2<5$.
The original constant $7H$ has therefore been preserved, rather than
replaced by an uncontrolled constant.

\subsection{Soundness of the definition}
Take a rational solution. Write $a_0=|a|$ and
$B=\max(|n|,|a|,|t|,|b|)$. The finite exclusion gives
$a_0>M>(20H)^2$. The unchanged scalar-height argument gives
\[
 h(\beta)\le H\sqrt{a_0}.
\]
For completeness about its use, its constant leaf is $1$ on the
all-left path of the binary norm tree, and the other six leaves are
the stated tests divided by $r_\beta$ to their assigned powers.
At every odd five-contact of $\beta$, the old norm and pole argument
gives order 48 for the parameter ratio, hence divisibility of the
integer label $a$ by the 48th power of that radical. The old
five-contact height theorem gives exactly the displayed estimate.

Set
\[
 \Delta=t^2-a(a-2)tb+(a-2)^2b^2-1\in\mathbb Z.
\]
If $\Delta\ne0$, the proxy equation and the precision-64 congruences
show that each of the nine radicals has 64th power dividing $\Delta$.
The source index-height inequality $h(x([m]P))\ge m^2-6$ and the
previous subsection give
\[
 B^2-6<7H|\Delta|^{3/8}<14HB^{3/2},
\]
where $|\Delta|<5B^4$ since $B\ge a_0>7$. Hence
$B<(14H+6)^2\le(20H)^2$, contradicting $B\ge a_0$.
Thus $\Delta=0$. The source's elementary Lucas-growth lemma, applied
to the nonzero label $b$, now gives
\[
 \log B>a_0-3.
\]

If $\beta\ne n$, let $N$ be the absolute numerator of $\beta-n$ in
lowest terms. The precision-16 comparisons show that every radical
has 16th power dividing $N$. As $N\le2e^{h(\beta)}B$, one obtains
\[
 B^2<7HN^{3/2}+6,
 \qquad \log B<2\log(32H)+3h(\beta).
\]
The bounds $a_0>(20H)^2$ and $H\ge1$ imply both
$h(\beta)<a_0/20$ and $\log(32H)<a_0/20$. Consequently
$\log B<a_0/4<a_0-3$, a contradiction. Therefore $\beta=n\in\mathbb Z$.
Since $\beta=X^2+2$ and $X\in\mathbb Q$, reduction of the denominator
of $X$ gives $X\in\mathbb Z$.

\subsection{Completeness and compatible simultaneous orders}

Fix an integer $X$ and put $n=X^2+2$. Use the source's scalar-height
completeness construction: choose $J_0$ in sufficiently deep formal
subgroups at the finitely many relevant primes, and then choose a
positive label $a$ sufficiently divisible relative to $J_0$. There
are arbitrarily large such labels, so require $a>M$. The source's
Lucas lemma supplies positive integers $t,b$ satisfying $\Delta=0$.
Set the proxies to $(n,a,t,b)$. This fixes the four root points, their
eight abscissa components and all positive square scales $A_r$.
The scalar norm root is a unit at every prime dividing $r_\beta$;
as $|r_\beta|>1$, it is nonzero. Thus the additional scalar guard
is satisfied and the original scalar pole witnesses exist.

Temporarily prescribe $H=[h]P$ and $U_i=[hn_i]P$. Each comparison
test is now a rational function of the point $[h]P$, obtained using
fixed multiplications and translations. At $h=64$ all its denominators
are nonzero: the relevant integer multiples are nonzero, $64\ne\pm4$,
and the source reciprocal logarithms never vanish there. Continuity
in the finitely many groups $E(K_{\mathfrak p})$, for
$\mathfrak p\mid2\cdot3\cdot17$, gives a positive integer $L$,
divisible by 64, and integers $c_p\ge0$ such that
\[
 h\equiv64\pmod L
 \quad\Longrightarrow\quad
 v_{\mathfrak p}(D_j)\ge-c_p
 \quad(\mathfrak p\mid p,\ p=2,3,17).
\]
One obtains $L$ by taking a common multiple of the indices of small
open formal subgroups contained in the required neighborhoods.

There is also a real bound $C$, independent of all future progression
moduli, attainable by arbitrarily large positive $h$ in every
nonconstant progression. Here is the compact-group argument. The
closure $G$ of $\{([h]P,[h]P^\sigma):h\in\mathbb Z\}$ in the
product of the real elliptic groups is a compact group with finitely
many components. Each component projects onto the identity circle,
up to translation, in each real elliptic factor: an infinite cyclic
projection is infinite, and its connected closure is the circle.
The finitely many poles of the required rational functions therefore
remove a set of Haar measure zero from each component. Choose a
neighborhood avoiding them in each component and bound all tests
there by one $C$. Multiplication by any positive integer maps $G^0$
onto itself, so the closure of any nonconstant progression contains
one component. Its positive tail is dense in that closure, proving
the assertion. This is the argument in Sun's Lemma~8.1; in particular
the bound is chosen before the prime pattern and its order moduli.

Choose a positive integer $D_0$ for which $D_0^2\theta/A_r$ is integral
for every contact $\theta$ and every $r$, and put
\[
 v=D_0^2w^2,\qquad \lambda=u/v.
\]
The contact positions are $u-(\theta/A_r)v$, a finite collection of
forms $u+c_jw^2$ with integer $c_j$. Delete repeated positions;
within each fixed $r$ its five positions are distinct. The zero
contact gives the common position $u$. Enlarge the prescribed lower
bound for the primes beyond $76176$ and beyond all prime factors of
the fixed coefficients and scales.

Invoke Sun's Theorem~4.1 with arbitrary prescribed divisibility of $w$,
arbitrarily large $u/w^2$, and the odd squarefree avoidance modulus
containing every odd prime factor of $3\cdot17L$ other than seven.
It supplies positive prime values at all positions, Gaussian
generators $\pi\equiv9\pmod{16}$, $\pi\equiv1+3i\pmod7$, and
$\pi\equiv i$ modulo that avoidance modulus. Its two reduction orders
are $d_+=N(\pi+1)$ and $d_-=N(\pi-1)$, with
\[
 v_2(d_+)=2,\qquad v_2(d_-)=6,\qquad
 \gcd(d_-(p),d_+(q))=4
\]
for all selected primes $p,q$. Both orders are prime to the odd part
of $L$: the prime seven follows separately from
$d_+\equiv6$, $d_-\equiv2\pmod7$.

Put $U=\operatorname{lcm}d_-$ and $V=\operatorname{lcm}d_+$.
The three congruences
\[
 h\equiv0\pmod U,\qquad h\equiv4\pmod V,\qquad
 h\equiv64\pmod L
\]
are compatible. Their pairwise common odd factors have been removed,
and the two-parts are compatible because
$v_2(U)=6$, $v_2(V)=2$, $64\mid L$, and $4\mid60$.
The generalized Chinese remainder theorem supplies a nonconstant
progression. Choose a sufficiently large positive $h$ in it attaining
the real bound $C$. This avoids $0,4,-4$ and every affine-chart
exception. At a selected prime choose the place with reduction order
$d_-$. Then $H,U_i$ are formal there and $V_0,V_i$ are formal at the
conjugate place of order $d_+$. Their exact two logarithm ratios equal
$n_i$, so all eighteen tests have their required local orders.

\subsection{Completeness of the nine pole occurrences}
The determinant of the contact coordinate change is a unit outside
the fixed exceptions. More explicitly, for each $r$,
\[
 R_r=\frac{\prod_{\theta\in\Theta_0}(u-(\theta/A_r)v)}
             {v^5(t_r-1)^6}.
\]
The denominator is a rational square, since $v$ is a square.
At its five distinct contact primes $R_r$ has valuation one, and
at every other prime it has even valuation. The contact primes do
not divide $v$; otherwise two distinct prime positions would share
that prime. They do not divide the denominator of a fixed scale or
a fixed difference, by the chosen lower bound. Thus $t_r-1$ is a
unit at these primes.

At either place above a contact prime, each factor of $W_r$ has
valuation zero or negative. Its inverse has residue one or zero,
respectively. At the successful place all factors have residue one.
Therefore $T_r$ has residue one or two and is a unit there. It follows
that $B_r=R_r/T_r^2$ has positive odd valuation at its contact primes
and even valuation at all other primes. In particular its reduced
denominator is a square.

The imposed Gaussian congruences give each contact prime
$p\equiv1\pmod8$ and $p\equiv1\pmod{3\cdot17}$. Such a prime is
square in each of $\mathbb Q_2,\mathbb Q_3,\mathbb Q_{17}$.
The displayed product formula for $R_r$, and then division by $T_r^2$,
shows that $B_r$ has the same local square properties. This is the
reason a single fixed multiplier $17$ suffices for all nine blocks.

It remains to arrange their quantitative preparation-zero conditions.
The real bound $C$ was uniform in the later congruence moduli, so for
$\lambda\to+\infty$ the fixed scales give, uniformly in the chosen $h$,
\[
 R_r\sim A_r^{-6}\lambda^{-1},\quad
 W_r=1+O(\lambda^{-3}),\quad T_r=2+O(\lambda^{-3}),\quad
 B_r\sim(4A_r^6\lambda)^{-1}.
\]
For sufficiently large $\lambda$ all factors at both real embeddings
are positive and $0<B_r<17$. At a small prime put
$\nu_p=v_p(v)=-v_p(\lambda)$ and $\eta_{p,r}=v_p(A_r)$.
For sufficiently large $\nu_p$, direct valuation calculation gives
\[
 v_p(R_r)=\nu_p-6\eta_{p,r},\qquad
 v_{\mathfrak p}(\widetilde D_{r,j})
 \ge-c_p+(m_j+1)(\nu_p-\eta_{p,r}).
\]
The nonconstant term of its factor consequently has valuation at least
\[
 3\nu_p+(10m_j-8)\eta_{p,r}-2c_p.
\]
These bounds tend to infinity with $\nu_p$. Thus $W_r$ and its inverse
are arbitrarily close to one at each place above $p$. Integrality of
trace, also at the ramified prime two, gives
$v_p(T_r)=v_p(2)$ once the error is sufficiently small. Hence
\[
 v_p(B_r)=\nu_p-6\eta_{p,r}-2v_p(2).
\]
Arbitrarily high prescribed divisibility of $w$ makes these valuations
meet $v_2(B_r)\ge3$, $v_3(B_r)\ge1$, $v_{17}(B_r)\ge2$ for every $r$.
All such denominator requirements and the real lower bound were chosen
before the application of the prime-pattern theorem. The restricted
pole-completeness lemma now supplies all nine structural pairs and
their square roots. This completes the rational solution for every
integral input.

\begin{theorem}[Conditional forty-six-variable construction]
Assume the arithmetic inputs explicitly stated at the beginning of
this section, including Sun's simultaneous prime-pattern theorem.
The rational system specified above has input projection precisely
$\mathbb Z$. It gives, by the finite-algebra construction, an
effectively specified polynomial over $\mathbb Z$ with at most
46 total variables, including its ordinary rational input. With the
three-witness semilocal anchor in place of the indexed finite exclusion,
the same conclusion holds with at most 48 total variables.
\end{theorem}
\begin{proof}
Soundness and completeness were proved in the preceding subsections.
The point substitution is exact by the single-chart lemma. On each
finite branch all denominators are guarded, all non-base choices lie
in the stated finite triangular algebra, and its zero-radicand and
repeated-root specializations have the prescribed finite treatment.
There are 45 retained coordinates and one final coordinate, or two
additional retained anchor coordinates in the semilocal version.
Clearing fixed rational coefficients gives integer coefficients.
The finite anchor is a specified indexed polynomial, not a claim that
its astronomically many factors have been expanded.
\end{proof}

\subsection{Preserving the scalar certificate when its point is shared}
\label{subsec:independent-height-sharing}

We verify that the scalar-height point $J_0$ can be identified with the
common comparison point $H$.  This is an arithmetic modification of the
system, with the same hypotheses as the shared-point construction.  The
scalar pole witnesses must be chosen after the final common index has
been selected.

\begin{proposition}
Under the stated source arithmetic inputs and Sun's simultaneous
prime-pattern theorem, imposing $J_0=H$ preserves the input projection
of the shared-point construction.  Consequently the indexed-exclusion
version uses at most $44$ total coordinates after finite-algebra
compression, and the semilocal-anchor version uses at most $46$.
\end{proposition}
\begin{proof}
Soundness follows from the construction with separate points by assigning
the same point to both occurrences.  For completeness fix an integer input,
write $r=r_\beta$, and choose the signed squarefree correction $c$ so that
$cr>0$ and $cr$ is square in $\mathbb Q_2,\mathbb Q_3,\mathbb Q_{17}$.
The number $cr$ is a nonzero integer.  Let $\mathcal S$ be its finite prime
support.

Choose a positive integer $h_0$ divisible by $64$ such that
$J_*= [h_0]P$ is sufficiently deep in the formal subgroups at all places
above $\mathcal S$.  Next choose the positive label $a$ of $S$ sufficiently
divisible relative to $h_0$, and sufficiently large for the size anchor.
The original scalar-completeness argument permits these successive choices.
It gives the following explicit conditions on the six scalar tests,
normalized by their assigned powers of $r$: every normalized test is
$p$-integral for $p\mid r$, and every nonconstant normalized test is divisible
by $p$ for $p\mid c$ with $p\nmid r$.  The denominators used in those tests
are nonzero.  The source Lucas construction now fixes the two other root
labels and all root points, and all four proxies are their integer labels.

These scalar conditions persist when $J_*$ is replaced by $[h]P$ in a
sufficiently small neighborhood at every place above $\mathcal S$.
Indeed each test is a rational function defined at $J_*$; the coefficient
maps from $K$ to its two rational components are continuous in the product
of the completions above $p$.  An integral lower valuation bound defines
an open condition at a point satisfying it.  Simultaneously impose the
uniform lower bounds for the eighteen comparison tests at the places
above $2,3,17$, using continuity at the nonzero indices $h_0$ and $h_0-4$.
Their denominators are nonzero by the reciprocal-logarithm lemma and
$64\mid h_0$, $h_0>0$.

The resulting finitely many local neighborhoods are preserved by one
integer progression
\[
                       h\equiv h_0\pmod L,
                       \qquad 64\mid L.
\]
To see this, put a sufficiently small open formal subgroup inside each
translated neighborhood, take a common multiple of its finite group
indices, and further multiply by $64$.

Use the compact real-group argument to obtain bounds for the comparison
tests uniform in every future progression modulus.  Choose the fixed
denominator divisibilities and real lower bound for the common slope as
in the separate-point construction.  In the simultaneous prime theorem,
include all odd prime factors of $3\cdot17L$, except seven, in its avoidance
modulus.  Its Gaussian class modulo seven already excludes seven from
both reduction orders.  With $U=\operatorname{lcm}d_-$ and
$V=\operatorname{lcm}d_+$, the hypotheses therefore give
\[
 \gcd(U,V)=4,\qquad \gcd(U,L)=64,\qquad \gcd(V,L)=4.
\]
The three congruences
\[
 h\equiv0\pmod U,\qquad h\equiv4\pmod V,
 \qquad h\equiv h_0\pmod L
\]
are compatible: $64\mid h_0$ and $4\mid h_0-4$ supply the two remaining
conditions.  Choose a sufficiently large positive $h$ in this progression
attaining the uniform real bounds, and put $J_0=H=[h]P$.
The comparison witnesses and the nine preparation-zero pole blocks are
then constructed exactly as before.

It remains to choose the scalar pole witnesses for this final $h$.
Let $W$ be its scalar norm root.  At a prime dividing $r$, the normalized
leaves are integral and the constant leaf $1$ on the all-left path makes
$W$ a unit.  At a prime dividing $c$ but not $r$, the nonconstant leaves
are divisible by that prime, so the same tree again has unit root.
These conditions hold because the progression preserved them.  Since
$|r|>1$, at least one such prime exists and $W\ne0$.

Thus $b_*=crW^2$ is positive.  At primes in $\mathcal S$ its valuation is
the nonnegative valuation of the integer $cr$; at all other primes its
valuation is even.  Its reduced denominator is therefore a square.
Its square classes at $2,3,17$ equal those of $cr$, so it is locally square
at all three places.  The original scalar pole-completeness theorem now
supplies its preparation and curve witnesses.  No real bound for $W$,
and no new rational coordinate, was required.

Sharing removes one point occurrence, hence two retained abscissa
coordinates and one quartic algebra generator.  The retained count becomes
$43$ (or $45$ with the semilocal anchor); the finite-algebra construction
adds one coordinate.  This proves the stated counts.
\end{proof}
